\newpage
\section{Extremwertaufgaben}
\label{4.5}

\begin{definition}{lokale Extrema}
    Sei $D\subset\mathbb{R}^n$ offen, $f: D \to \mathbb{R}$.

    \medskip

    Ein Punkt$x\in D$ heißt \emph{lokales Minimum (Maximum)} von $f$, falls eine Umgebung$B_{\delta}(x)\subset\mathbb{R}^n$ existiert mit
    $$
    f(x)\le f(y),\quad \forall y\in B_{\delta}(x)\cap D
    $$
    bzw.
    $$
    f(x)\ge f(y),\quad \forall y\in B_{\delta}(x)\cap D.
    $$

    \medskip

    Falls
    $$
    f(x)< f(y),\quad \forall y\in B_{\delta}(x)\cap D\setminus\{x\}
    $$
    (bzw. $f(x)>f(y)$),
    dann heißt$x$ \emph{striktes lokales Minimum (Maximum)}.
\end{definition}

\begin{satz}{Notwendige Bedingung für lokales Extremum}
    Sei $D\subset\mathbb{R}^n$ offen, $f\colon D\to\mathbb{R}$ stetig differenzierbar und $x\in D$ ein lokales Extremum von $f$.

    Dann gilt:
    $$
    \nabla f(x)=0.
    $$
\end{satz}

\begin{satz}{Hinreichende Bedingung für lokales Extremum}
    Sei $D\subset\mathbb{R}^n$ offen, $f\in C^2(D,\mathbb{R})$ und $x\in D$ mit $\nabla f(x)=0$. Dann gilt:
    \begin{enumerate}
        \item$H_f(x)$ positiv definit$\Longrightarrow$$x$ striktes lokales Minimum von $f$.
        \item$H_f(x)$ negativ definit$\Longrightarrow$$x$ striktes lokales Maximum von $f$.
        \item$H_f(x)$ indefinit$\Longrightarrow$$x$ kein lokales Extremum.
    \end{enumerate}
\end{satz}